\section{Implementation and verification status}
\label{sec:implementation}

This section states precisely what is implemented, what is verified, and
what is not.  We hold this to be part of the language's design: because R6
requires promises and caveats to travel with programs, the same discipline
applies to the paper.

\subsection{Implementation}
\label{sec:implementation-status}

\pqlang{} is a pure-Python package (version 0.8.0; Python $\ge 3.11$) with
zero third-party runtime dependencies.  Source is layered as
\texttt{infrastructure} (IR, builder, validation, serialization, math
frontend, backends), \texttt{algorithms} (the primitive and solver
libraries), and \texttt{applications} (QFVM, QHAM mathematics, the catalog).
The RIR specification (currently version 0.3), its JSON Schema, the
serializer, and the semantic tests are maintained in lockstep---a norm from
the project's development conventions---and readers accept prior versions,
so stored programs remain loadable across format revisions; a condensed
form of the specification appears in Appendix~\ref{app:spec}.  Native backends
(PySparQ; the uniqc OriginIR simulator) are optional environments pinned in
a reviewed revision ledger rather than installed dependencies; generation
and serialization work without them, and integration tests run in whichever
interpreter provides them, without mocks or substitutes.

\subsection{A graded evidence ladder}
\label{sec:evidence-ladder}

Verification claims are organized as levels, defined in the project's
validation plan and enforced by its test suite:

\begin{itemize}
\item \textbf{L0 -- static}: schema validation of every serialized format;
  ruff-clean sources; documentation built with warnings as errors, with
  doctests executing every tutorial snippet quoted in this paper.
\item \textbf{L1 -- structural validity}: IR validation (typing, alias and
  control protection, call-graph acyclicity, capability demands),
  serialization round-trips (byte-identical, open modules included), and
  contract checking of algorithm inputs.
\item \textbf{L2 -- classical witnesses}: small-scale reference computations
  attached to input models (e.g.\ Fokker--Planck stationary distributions,
  QHAM linearization identities, walk hitting-time estimates) checked in
  tests.
\item \textbf{L3 -- catalog}: every catalog case exports, parses on real
  parsers, and satisfies its declared binding plan.
\item \textbf{L4 -- real backends}: full-amplitude agreement between the
  reference simulator and both native backends for the 22-case gallery,
  math-function compilations, and gate-vs-QRAM oracle realizations, to
  $10^{-10}$; no simulator substitutes.
\item \textbf{L5 -- documentation}: doctest builders over the manual.
\end{itemize}

What the ladder does \emph{not} by itself include is numerical certification
of the scientific solver stack: solution accuracy against continuum
references, success-channel probabilities, and convergence in method
parameters.  Those are properties of algorithms, not of circuits.  We
therefore ran a dedicated numerical-verification campaign on the real
backends, reported in \secref{sec:numverify}; its outcomes are folded into
the last column of Table~\ref{tab:status-matrix}, and the remaining
open items are tracked per method as explicit pending items attached to
generated modules (R6).

\begin{table}[t]
  \centering
  \caption{Verification status by component.  ``Structure'' = L1 validity,
  serialization, contract checks; ``witness'' = L2/L3 classical witnesses
  and catalog reproducibility; ``native circuits'' = L4 real-backend
  amplitude agreement of assembled circuits; ``numerical certification'' =
  solution-accuracy and success-channel checks against mathematical
  references in the campaign of \secref{sec:numverify} (register-exact
  scale; scaling studies remain open).  A dash means \emph{not yet
  performed}, not impossible.}
  \label{tab:status-matrix}
  \small
  \setlength{\tabcolsep}{4pt}\begin{tabular}{@{}p{4.3cm}cccc@{}}
    \toprule
    \tabhead{Component} & \tabhead{structure} & \tabhead{witness} & \tabhead{native circuits} & \tabhead{num.\ cert.} \\
    \midrule
    RIR core, builder, serialization & \ok & \ok & \ok & --- \\
    Oracle lifecycle (\texttt{bind}, capabilities) & \ok & \ok & \ok & --- \\
    QSP phase synthesis / QSVT transforms & \ok & \ok & \ok & \ok$^{a}$ \\
    Data loading (QROM, QRAM, alias sampling) & \ok & \ok & \ok & \ok \\
    Reversible arithmetic \& math functions & \ok & \ok & \ok & \ok \\
    Estimation suite (QPE, AE, Hadamard, swap) & \ok & \ok & \ok & \ok \\
    Hamiltonian simulation (Trotter, Taylor) & \ok & \ok & \ok & \ok \\
    QLSS kernels (Costa walk, CKS Chebyshev) & \ok & \ok & \ok$^{b}$ & \ok \\
    ODE families (LCHS, Schr., CBMD, Carleman, Euler--QLSS, Taylor) & \ok & \ok & \ok$^{b}$ & partial$^{c}$ \\
    PDE layer (heat, Fokker--Planck, QHAM) & \ok & \ok & \ok$^{b}$ & \ok \\
    QFVM / Euler pipeline & \ok & \ok$^{d}$ & \ok$^{b}$ & partial$^{e}$ \\
    Scientific suite (Gibbs, QPCA, SDP, DQI, \dots) & \ok & \ok & \ok & \ok \\
    Resource estimator (strict gate set) & \ok & \ok$^{f}$ & --- & --- \\
    \bottomrule
    \multicolumn{5}{@{}p{14.3cm}@{}}{\scriptsize $^a$QSP synthesis self-verifies polynomial round-trips; transforms are certified end-to-end in \secref{sec:numverify}.  $^b$assembled-circuit agreement on real backends.  $^c$all six families certified against truncated-oracle references at small scale; the Schr\"odingerization momentum-sign defect found by an earlier campaign round is repaired and pinned by regression cases (\secref{sec:numverify}); success channels and large-scale accuracy remain pending.  $^d$classical flux and flow-state witnesses.  $^e$fixed-point physics certified against an independent fixed-point semantics at \texttt{FixedFormat(5,2)}; generation-time constant quantization dominates the method error in coarser formats.  $^f$hand-computed counts and atom-by-atom cross-checks against the strict netlist; scaling expectations validated in Appendix~\ref{sec:resources}.}
  \end{tabular}
\end{table}

\subsection{Numerical verification on real backends}
\label{sec:numverify}

We close the certification column at the scale where exact classical
references exist, using a dedicated numerical campaign rather than unit
tests.  Two pieces of backend infrastructure were built for it and now ship
with the respective simulators.  First, PySparQ executes RIR natively:
\texttt{pysparq.rir} parses the versioned JSON encoding and interprets it
directly on sparse states, expanding the module graph (calls,
\texttt{Repeat}, \texttt{Control}, \texttt{Adjoint}) at interpretation time
and mapping register-aligned arithmetic onto native integer operators---so
the simulator, not only the compiler, is now an independent RIR
implementation.  Second, UnifiedQuantum circuits expose
\texttt{Circuit.to\_matrix()}, the full unitary of a parsed OriginIR-ext
listing including modular \texttt{DEF} structure, which lets us compare
compiled arithmetic against classical permutation matrices elementwise.
Together with the pre-existing reference executor and the event-level
PySparQ adapter, every program in the campaign can be executed by up to
four independently written implementations, cross-checked wherever budgets
allow.  Differential testing of this kind caught a genuine wraparound
defect in the new interpreter (addition on a register slice wrapping at the
register width instead of the view width), which is fixed and covered by
regression tests on both sides.

The campaign (Table~\ref{tab:verify-matrix}) comprises 475 experiments in
thirteen groups, driven by \texttt{tools/run\_verification.py} and archived
as machine-readable JSON (\texttt{out/verification/}).  Every experiment
compares against a correctness oracle that is independent of the
implementation---closed forms, \texttt{numpy}/\texttt{scipy} references, or
an independently written fixed-point semantics---and most sweeps evaluate
the entire input domain in a single superposed execution instead of
sampling basis states.  Where an algorithm combines an exact construction
with a mathematical truncation (Chebyshev degrees, quadrature cutoffs,
Carleman tails), \emph{implementation error} (against the truncated oracle)
and \emph{method error} (of the truncation itself) are reported separately;
domain flags, padding structure, and workspace cleanliness are checked
branch by branch.

\begin{table*}[t]
  \centering
  \caption{Numerical-verification campaign: thirteen groups, 475 experiments,
  all passing their stated criteria on real backends.  Paths: O =
  OriginIR-ext + UnifiedQuantum full amplitudes (incl.\
  \texttt{to\_matrix}); P = PySparQ RIR sparse-state execution; X =
  cross-checked against the reference executor and the event adapter.
  Headlines quote the worst-case metric over the group's instances.}
  \label{tab:verify-matrix}
  \small
  \setlength{\tabcolsep}{4pt}\begin{tabular}{@{}p{3.6cm}ccp{8.6cm}@{}}
    \toprule
    \tabhead{Group} & \tabhead{cases} & \tabhead{paths} & \tabhead{headline metrics} \\
    \midrule
    Reversible arithmetic \& fixed point & 20 & O, P, X & adders permutation-exact to 64 bit; compiled polynomials $\le 1$ quantum; $\sin$ impl.\ $0.055$ / method $0.045$ (deg.\ 3) \\
    Fourier \& unitary transforms & 41 & O, P, X & QFT vs DFT $3.8{\times}10^{-15}$ ($n{=}4$); Fourier addition $1.0{\times}10^{-15}$; OAA amplifies $\sin\theta \to \sin 3\theta$ exactly after repair ($0.5 \to 1.0$; unitary $\le 4.6{\times}10^{-16}$) \\
    State preparation \& data loading & 37 & O, P, X & fidelity $\ge 1 - 2{\times}10^{-16}$; QROM/select-swap $\le 8.4{\times}10^{-17}$; alias TVD within analytic bounds \\
    Oracle algorithms & 30 & O, P, X & BV recovery $= 1$ ($w = 3$--$8$, both paths); DJ/Simon exact; XOR database exhaustive $5.6{\times}10^{-17}$ \\
    Search \& quantum walks & 16 & O, P, X & Grover curve $2.3{\times}10^{-15}$ vs $\sin^2(2k{+}1)\theta$; MNRS vs independent reference $1.2{\times}10^{-15}$ \\
    Block encodings & 73 & O, P, X & diagonal/sparse/LCU blocks $\le 1.1{\times}10^{-16}$; PySparQ BE cross-check $\le 6.7{\times}10^{-16}$; Taylor method trend $0.141 \to 7.8{\times}10^{-4}$ \\
    Hamiltonian simulation \& QSVT & 27 & O, P, X & Trotter order $1.008$/$1.013$ measured (theory $1$); QSVT HamSim total $\le 4.5{\times}10^{-4}$; inversion within prescribed error \\
    Estimation suite & 41 & O, P, X & QPE Dirichlet TVD $\le 7.3{\times}10^{-16}$; QAE within the Brassard bound; Jordan gradient exact on linear oracles \\
    ODE solvers & 27 & O, P, X & LCHS impl.\ $\le 1.7{\times}10^{-16}$ over five input models; CBMD direction error $3.0{\times}10^{-3}$ vs LCHS $4.5{\times}10^{-2}$; Carleman tail $9.3{\times}10^{-3} \to 1.1{\times}10^{-4}$; Schr.\ forward-flow recovery after sign repair (grid $1.3{\times}10^{-16}$, Taylor remainder $3.3{\times}10^{-2}$) \\
    QHAM \& QFVM & 15 & P, X & generator matrices $\le 1.5{\times}10^{-17}$; input-model agreement; Roe face bit-exact vs independent fixed-point semantics \\
    Math-function gallery & 28 & P, X & five gallery functions + Roe face bit-exact vs independent fixed-point semantics; domain flags exact \\
    Density, QPCA, SDP, DQI, variational, QEC & 24 & O, P, X & LMR trace distance with first-order scaling ratio $\approx 0.51$; DQI planted instance TVD $4{\times}10^{-16}$; repetition codes exact \\
    Number theory, Heinrich, SDE, QLSS, catalog & 96 & O, P, X & modular multiplication $\le 8.3{\times}10^{-17}$; order finding vs Dirichlet $\le 8.6{\times}10^{-16}$; CKS kernel $\le 4.9{\times}10^{-16}$; SDE Euler order $1.0001$; catalog 33 cases $\le 1.4{\times}10^{-16}$ \\
    \bottomrule
  \end{tabular}
\end{table*}

Three findings are worth reporting explicitly, because they are exactly the
kind of evidence a language-level verification story should produce.
\emph{First, the campaign found real defects---and repaired them with
regression coverage.}  The oblivious-amplitude-amplification iterator
originally assembled $R\,U^\dagger R\,U$, realizing its documented operator
to $10^{-16}$, but that operator is not the amplifying sequence; the
repaired form $U\,[R\,U^\dagger R\,U]^{m}$ equals the standard three-query
sequence, satisfies the Chebyshev identity
$\Pi W \Pi = B(4B^\dagger B - 3I)$ on its signal block, and amplifies
$0.5 \to 1.0$ exactly, with the library assembly asserted equal to an
independently reassembled sequence ($\le 10^{-15}$) as a regression pin.
Schr\"odingerization likewise assembled its advertised generator
faithfully, yet recovered the \emph{time-reversed} flow to
$4.8{\times}10^{-16}$ under the library's positive-QFT convention; the
repair flips the momentum sign of the lifted generator
($K' = -P\otimes H_1 - I\otimes H_2$), restoring the forward flow to
$1.3{\times}10^{-16}$ under exact evolution, with the end-to-end residual
equal to the Taylor remainder of the Nyquist momentum mode
($3.3{\times}10^{-2}$ at degree four, decreasing with degree).  Both
defects are reported openly with their repairs, per R6.
\emph{Second, the dominant error in the fixed-point scientific kernels is
not circuit noise but generation-time constant quantization} (e.g.\
$\gamma - 1 = 0.4 \mapsto 0.25$ at six-bit formats): implementation error
against the compiled semantics stays at or below one quantum everywhere,
and the quantization contribution is the leading term of the method error.
\emph{Third, there is a measurable backend floor}: on deeply nested LCU
programs ($\sim 10^6$ expanded steps) the sparse-state backends prune junk
branches below $\sim 10^{-7}$, while the reference executor and the
OriginIR-ext path agree with the polynomial oracles at $10^{-15}$; all
physical post-selected blocks remain consistent across paths.

The scope of this certification is deliberately register-exact: instances
are small enough that amplitudes, blocks, and success probabilities can be
compared essentially to machine precision, and the reported method errors
characterize the mathematical truncations, not estimates of large-scale
behaviour.  Scaling the campaign is the subject of the next subsection.

\subsection{Resource estimation in brief}
\label{sec:resources-summary}

Beyond correctness, the same structure preservation supports compositional
resource estimation: \texttt{export\_strict} lowers a program to a
module-preserving Toffoli+Clifford+T+QRAM netlist and
\texttt{estimate\_resources} counts qubits, Toffoli, synthesis-adjusted
$T$, and QRAM queries per (module, control-count) pair, multiplying
through \texttt{Repeat} nodes without expansion.  We treat this as a
secondary capability of the architecture rather than a core contribution:
the strict-gate-set compiler, the scaling study against theoretical
expectations, the query-count view of QRAM, and the solver-stack tables
are reported in Appendix~\ref{sec:resources}.

\subsection{What remains reserved}
\label{sec:reserved}

The remaining axis of validation---and in our view the decisive one for a
scientific-computing language---is deliberately \emph{reserved} for future
work and summarized here only as a statement of scope:

\begin{itemize}
\item \textbf{Backend validation at scale.}  L4 evidence currently covers
  small instances where full amplitudes are checkable.  Large, symbolic-scale
  programs (big repetition counts, wide QRAM) need staged execution regimes
  and differential testing across the native backends; the cost-model guards
  and lazy expansion of \secref{sec:rir} are the enabling machinery, but the
  campaigns themselves are future work.
\item \textbf{Simulation campaigns for the solver stack at scale.}  The
  campaign of \secref{sec:numverify} certifies implementations against
  truncated-oracle references at register-exact scale.  What remains is the
  parameter-scaling axis---quadrature cutoffs for LCHS, recovery windows
  for Schr\"odingerization, contour plans for CBMD, Carleman cutoffs,
  Taylor degrees, and success channels as functions of problem size---and
  the repair of the two recorded defects.  The promise-provenance
  attributes are designed to be the bookkeeping substrate for exactly these
  studies.
\item \textbf{Resource estimation, physical layer.}  The logical level
  landed in Appendix~\ref{sec:resources}: a strict Toffoli+Clifford+T+QRAM
  compile and a compositional estimator derived from the module graph
  itself, in the spirit of automated application-level frameworks
  \cite{beverland2022assessing,harrigan2024qualtran,su2026resource} and of
  large-scale studies such as \cite{gidney2021factor,gidney2025rsa}.  What
  remains open is everything below that level: rotation synthesis as an
  implemented pass rather than an accounted model, depth scheduling,
  ancilla and distillation-factory budgets, routing, and surface-code cycle
  accounting---as well as flowing the measured numbers back into per-module
  promise attributes (R6) and repairing the lowering gaps the estimator
  exposed (QROM emission, constant-addition staircase).
\end{itemize}

\noindent The physical layer and the scaling axis are deliberately scoped
as future work rather than reported prematurely: attaching physical-cost
numbers to an uncertified-at-scale solver stack would violate the very
provenance discipline this paper argues for.
